\section{A rank-preserving normalized frame}
\label{sec:normalization}


The shortcut controls effective resistances. A second vector system
encodes the spanning-tree packing, so that signing can preserve both
properties simultaneously. The normalization uses the log-determinant
variational method for radial isotropic position
\cite[Section~2.2]{Barthe1998}. We prove the approximate positive-weight
form and its bit bounds directly, so no attainment assumption is needed.

\Needspace{4\baselineskip}
\begin{lemma}[A computable normalized frame]
\label{lem:ta-frame}
Suppose $B$ is the union of $q$ edge-disjoint spanning trees, where
$r/2\le q\le r$ and $r$ exceeds an absolute constant. There are vectors
$z_e\in\R^{n-1}$, obtained from faithful reduced incidence columns by one
common invertible linear map and positive individual scalings, such that
\begin{equation}
 \sum_{e\in B}z_ez_e^T=I,\qquad
                  \|z_e\|^2\le\frac{1+4/r}{q}.
 \label{eq:ta-frame}
\end{equation}
They admit rational approximations to Euclidean error $2^{-b}$ in time
polynomial in $n+|B|+r+b$. An exact specification uses polynomial-bit
positive rational edge weights followed by whitening; in particular,
every subset has exactly its original incidence rank.
\end{lemma}

\begin{proof}
Delete one row of the incidence matrix and write the resulting columns
as $c_e$. Put $M=|B|=q(n-1)$ and consider
\begin{equation}
 N(d)=\sum_{e\in B}\exp(d_e)c_ec_e^T,\qquad
 f(d)=\log\det N(d)-\frac1q\sum_{e\in B}d_e.
 \label{eq:ta-logdet}
\end{equation}
By Cauchy--Binet,
$\det N(d)=\sum_T\exp(\sum_{e\in T}d_e)$, where $T$ ranges over the
spanning trees of $B$: reduced incidence determinants of trees have
absolute value one, and other maximal minors vanish. Keeping the $q$
specified disjoint trees and applying the arithmetic--geometric mean
inequality gives
\[
 \det N(d)\ge q\exp\left(\frac1q\sum_{e\in B}d_e\right),
               \qquad f(d)\ge\log q.
\]
Also $f(0)\le(n-1)\log(2M)$, so the initial gap above the infimum is
bounded by $n\log(2M)$. Differentiation gives
\begin{equation}
 \partial_e f(d)=\exp(d_e)c_e^TN(d)^{-1}c_e-\frac1q.
 \label{eq:ta-gradient}
\end{equation}
The Hessian is the covariance matrix of the tree indicator under the
probability distribution proportional to
$\exp(\sum_{e\in T}d_e)$. Hence it is PSD with operator norm at most
$M$, since these indicators have squared norm $n-1\le M$.
Each leverage score in \Cref{eq:ta-gradient} lies in $[0,1]$; consequently
$\|\nabla f\|\le\sqrt M$.

Start at $d=0$. At each step compute a rational gradient estimate $v$
with $\|v-\nabla f(d)\|\le1/(10r^2)$. If $\|v\|\le2/r^2$, stop;
otherwise update
$d\leftarrow d-v/(2(M+1))$. The norm comparison can be made by squaring
rational quantities. The Hessian bound yields
\[
 f(d-\alpha v)\le f(d)-\alpha\nabla f(d)\cdot v
                         +\tfrac12M\alpha^2\|v\|^2,
                       \qquad \alpha=\frac1{2(M+1)}.
\]
At a nonstopping step the gradient error is at most $\|v\|/20$.
Thus the decrease is at least
$\alpha\|v\|^2/2\ge2\alpha/r^4$.
The gap bound permits at most
$O((M+1)r^4 n\log(2M))$ such steps. At termination
\begin{equation}
                       \|\nabla f(d)\|\le3/r^2.
 \label{eq:ta-gradient-stop}
\end{equation}

Here is an explicit precision justification for this descent. Its
iteration bound is prescribed in advance, and each coordinate changes
by at most two in one step. Thus $|d_e|\le B_0$ for a known polynomial
$B_0$ in $n+M+r$. On that entire box,
\[
 e^{-B_0}N(0)\preceq N(d)\preceq e^{B_0}N(0).
\]
The matrix $N(0)$ is positive definite, has norm at most $2M$, and has
positive integral determinant. Its smallest eigenvalue is therefore
at least $(2M)^{-(n-2)}$. The upper norm and inverse-norm bounds for
$N(d)$ have polynomial logarithms. Computing exponentials to
polynomially many bits, forming rational approximations to $N(d)$,
and solving rational linear systems therefore evaluates
\Cref{eq:ta-gradient} to the requested error in polynomial bit work.
For example, the identity
$\widetilde N^{-1}-N^{-1}=\widetilde N^{-1}(N-\widetilde N)N^{-1}$
controls the inverse error once
$\|\widetilde N-N\|\le\lambda_{\min}(N)/2$.
Gradient estimates may be rounded on one fixed dyadic grid $2^{-b_0}$
sufficiently fine for the stated error. Starting at zero, every iterate
then lies on the grid $[2(M+1)2^{b_0}]^{-1}\mathbb Z^M$.
The prescribed iteration and magnitude bounds therefore keep every
$d_e$ a polynomial-bit rational.

At the terminal $d$, choose positive rational weights $a_e$ satisfying
$|a_e/\exp(d_e)-1|\le1/(100r^2)$ and put
$N_a=\sum_e a_ec_ec_e^T$.
Compute its rational diagonal elimination $N_a=L D_aL^T$ with positive
diagonal $D_a$, set $F=D_a^{-1/2}L^{-1}$, and define
$z_e=\sqrt{a_e}Fc_e$. Then $FN_aF^T=I$. The relative matrix bounds on $N_a$ and
\Cref{eq:ta-gradient-stop} imply
\[
 \|z_e\|^2
 \le\frac{1+1/(100r^2)}{1-1/(100r^2)}
                          \left(\frac1q+\frac3{r^2}\right)
 \le\frac1q+\frac4{r^2}\le\frac{1+4/r}{q}.
\]
Their outer products sum exactly to $I$. All $a_e$ are strictly positive
and $F$ is invertible, proving rank preservation.

For the promised approximations, perform exact rational diagonal
elimination on the positive definite $N_a$ and then approximate the
positive square roots of its pivots. Every pivot and inverse pivot
has polynomial logarithmic bounds, either by the preceding spectral
bounds or by determinant bounds for rational principal submatrices.
Triangular solves and scalar square roots therefore approximate $F$
and all $z_e$ to $2^{-b}$ with polynomially many guard bits and
polynomial work. These approximations refer to the single exact frame
just defined; rounding does not redefine its subset ranks.
\end{proof}
